%-------------------------
% Resume in LaTeX (Chinese, Internship Version)
% Compile with: xelatex resume-zh-intern.tex
%-------------------------

\documentclass[letterpaper,11pt]{article}

\usepackage{latexsym}
\usepackage[empty]{fullpage}
\usepackage{titlesec}
\usepackage{marvosym}
\usepackage[usenames,dvipsnames]{color}
\usepackage{verbatim}
\usepackage{enumitem}
\usepackage[hidelinks]{hyperref}
\usepackage{fancyhdr}
\usepackage{tabularx}
\usepackage{fontspec}
\usepackage{xeCJK}
\setCJKmainfont{WenQuanYi Zen Hei}
\setmainfont{Latin Modern Roman}

\pagestyle{fancy}
\fancyhf{}
\fancyfoot{}
\renewcommand{\headrulewidth}{0pt}
\renewcommand{\footrulewidth}{0pt}

\addtolength{\oddsidemargin}{-0.5in}
\addtolength{\evensidemargin}{-0.5in}
\addtolength{\textwidth}{1in}
\addtolength{\topmargin}{-.9in}
\addtolength{\textheight}{1.8in}

\urlstyle{same}
\raggedbottom
\raggedright
\setlength{\tabcolsep}{0in}

\titleformat{\section}{
  \vspace{-4pt}\scshape\raggedright\large
}{}{0em}{}[\color{black}\titlerule \vspace{-5pt}]

\newcommand{\resumeItem}[1]{
  \item\small{
    {#1 \vspace{-2pt}}
  }
}

\newcommand{\resumeSubheading}[4]{
  \vspace{-2pt}\item
    \begin{tabular*}{0.97\textwidth}[t]{l@{\extracolsep{\fill}}r}
      \textbf{#1} & #2 \\
      \textit{\small#3} & \textit{\small #4} \\
    \end{tabular*}\vspace{-7pt}
}

\newcommand{\resumeProjectHeading}[2]{
    \item
    \begin{tabular*}{0.97\textwidth}{l@{\extracolsep{\fill}}r}
      \small#1 & #2 \\
    \end{tabular*}\vspace{-7pt}
}

\renewcommand\labelitemii{$\vcenter{\hbox{\tiny$\bullet$}}$}

\newcommand{\resumeSubHeadingListStart}{\begin{itemize}[leftmargin=0.0in, label={}]}
\newcommand{\resumeSubHeadingListEnd}{\end{itemize}}
\newcommand{\resumeItemListStart}{\begin{itemize}[leftmargin=*, itemsep=0pt, topsep=2pt]}
\newcommand{\resumeItemListEnd}{\end{itemize}\vspace{-5pt}}

%-------------------------------------------
\begin{document}

%----------HEADING----------
\begin{center}
    \textbf{\Huge \scshape 江嘉元} \\ \vspace{1pt}
    \small +886 928 530 700 $|$ \href{mailto:johnjia9491@gmail.com}{johnjia9491@gmail.com} $|$
    \href{https://github.com/Jcsk7049}{github.com/Jcsk7049} $|$
    \href{https://linkedin.com/in/\%E5\%98\%89\%E5\%85\%83-\%E6\%B1\%9F-662764344}{LinkedIn} $|$
    台北市，台灣
\end{center}

%----------OBJECTIVE----------
\vspace{-4pt}
\small{\textbf{求職目標：} 電機系大三在學，尋找嵌入式韌體／FAE／軟硬整合實習。從機構、CNC 加工、PCB、嵌入式韌體到資料分析皆親手做過。暑假可全職實習，學期間可兼職。}
\vspace{2pt}

%-----------EDUCATION-----------
\section{學歷}
  \resumeSubHeadingListStart
    \resumeSubheading
      {元智大學}{桃園市}
      {電機工程學系 學士}{2024.09 -- 2027.09（預計）}
    \resumeSubheading
      {南港高級工業職業學校}{台北市}
      {電子科（機電整合）}{2020.09 -- 2023.06}
  \resumeSubHeadingListEnd

%-----------PROJECTS-----------
\section{專案}
    \resumeSubHeadingListStart
      \resumeProjectHeading
          {\textbf{開源 QMK x STM32 數字鍵盤} $|$ \emph{C, ChibiOS HAL, EasyEDA}}{2026}
          \resumeItemListStart
            \resumeItem{以 EasyEDA Pro 設計 17 鍵 PCB，手焊組裝並透過 DFU 燒錄 STM32F072；逆向工程 ChibiOS HAL 三件組（chconf/halconf/mcuconf），解決開源 PCB 與 QMK 韌體間的編譯衝突與時序問題。}
            \resumeItem{達成穩定 USB HID 識別（按鍵延遲 <5ms）、RGB Matrix 正常運作，並支援 VIA/VIAL 即時改鍵、無需重新編譯。}
          \resumeItemListEnd
      \resumeProjectHeading
          {\textbf{ICU VAP 早期預警模型} $|$ \emph{TensorFlow/Keras, LSTM, Integrated Gradients}}{2025 -- 至今}
          \resumeItemListStart
            \resumeItem{單中心研究（n=109）；找出 PREDICT 基準採用 window-level 交叉驗證造成的資料洩漏（同一病人的時間窗橫跨訓練/測試集），改設計 stay-level 5-fold 分層交叉驗證，反映對未見病人的真實泛化能力。}
            \resumeItem{以 Integrated Gradients 將模型精簡至 4 個非侵入性特徵；最終 LSTM 在 6--72 小時窗口維持 0.98--0.99 AUROC。論文已投稿 IEEE GCCE 2026。}
          \resumeItemListEnd
      \resumeProjectHeading
          {\textbf{Job Radar} $|$ \emph{Python, Playwright, GitHub Actions}}{2026 -- 至今}
          \resumeItemListStart
            \resumeItem{自行打造職缺聚合器，每日透過 GitHub Actions 排程爬取 104／Cake／LinkedIn（約 1,300 筆/天），依個人條件評分排序並產出靜態儀表板；無後端、電腦關機也照常更新。}
          \resumeItemListEnd
      \resumeProjectHeading
          {\textbf{Swerve Drive 全向輪底盤（FRC）} $|$ \emph{LabVIEW, CNC, PID Control}}{2020 -- 2023}
          \resumeItemListStart
            \resumeItem{從零開發三代全向輪底盤（運動學演算法、齒輪傳動、CNC/線切割加工、LabVIEW PID 控制），獲 2022 FRC 台灣鴻海區域賽控制創新獎。}
          \resumeItemListEnd
    \resumeSubHeadingListEnd

%-----------EXPERIENCE-----------
\section{經歷}
  \resumeSubHeadingListStart
    \resumeSubheading
      {專題研究生}{2025.05 -- 至今}
      {元智大學 林書彥教授實驗室}{桃園市}
      \resumeItemListStart
        \resumeItem{主導上述 ICU VAP 早期預警專題；負責從資料清理、交叉驗證設計到模型訓練與特徵歸因的完整流程。}
      \resumeItemListEnd
    \resumeSubheading
      {工程顧問與系統整合}{2023 -- 至今}
      {FRC Team 7645}{台北市}
      \resumeItemListStart
        \resumeItem{負責機器人系統架構：多感測器融合、PID 閉迴路馬達控制調校、軟硬體介面定義；訓練新成員以根因分析取代試錯法除錯。}
      \resumeItemListEnd
  \resumeSubHeadingListEnd

%-----------TECHNICAL SKILLS-----------
\section{技術能力}
 \begin{itemize}[leftmargin=0.0in, label={}]
    \small{\item{
     \textbf{程式語言與韌體}{: C/C++ (STM32/Cortex-M), Python, QMK/ChibiOS HAL, LabVIEW, Arduino/ESP32, UART/SPI/I2C} \\
     \textbf{EDA / 硬體設計}{: Altium Designer, KiCAD, OrCAD, EasyEDA Pro} \\
     \textbf{機器學習}{: PyTorch, TensorFlow/Keras, LSTM/時間序列, Integrated Gradients/XAI, TinyML} \\
     \textbf{製造}{: CNC 銑床, Mastercam, 3D 列印, 雷射切割, 電焊, AutoCAD, Fusion 360} \\
    }}
 \end{itemize}

%-----------AWARDS & CERTIFICATIONS-----------
\section{獎項與證照}
 \begin{itemize}[leftmargin=0.0in, label={}]
    \small{\item{
     2022 FRC 台灣鴻海區域賽 -- 決賽聯盟與控制創新獎 $\cdot$
     第 52 屆全國技能競賽 -- 團隊專案類代表（2022） $\cdot$
     2020 FRC CIT 5G 區域賽 -- 工程設計卓越獎
    }}
 \end{itemize}

\end{document}
